Iterative ``ask an AI to review, then revise, then repeat'' loops are increasingly used in automated editing and agentic writing pipelines, on the implicit assumption that they eventually \emph{stabilize} at a polished, fixed document. We test this assumption directly. We formalize the loop as a discrete dynamical system $V_{t+1} = \Revise(V_t, \Review(V_t))$ and ask whether it reaches a fixed point where further iterations produce no changes. Across 12 scientific abstracts, 8 iterations, and 8 experimental conditions, we find that the decoupled two-stage review$\to$revise loop \emph{does not} converge: the per-step normalized edit distance plateaus at a large residual floor (mean final $\Delta \approx 0.38$ versus an approximate-fixed-point threshold of $0.02$; one-sample $t=11.3$, $p=2\times10^{-7}$), and $0\%$ of trajectories reach a fixed point. The cause is the explicit review stage. A single-stage self-refinement baseline on the \emph{same} documents does converge (mean final $\Delta \approx 0.05$, $42\%$ reach an approximate fixed point, exponential-relaxation fit $R^2=0.98$), replicating prior work; inserting a freshly generated peer review at every step turns this convergent relaxation into a non-converging expansion process. A never-satisfied reviewer always emits new demands, which the reviser satisfies by adding text: documents inflate $\sim 7\times$, drift semantically from the original (cosine $0.97\to 0.77$), and lose factual fidelity (meaning preservation $10\to 6.2$ out of 10)—yet isolated quality scores barely move, so the loop \emph{feels} like it is improving while diverging. Running such loops ``to stability'' is not a reliable stopping strategy.
